\documentclass[12pt]{article}
\usepackage[utf8]{inputenc}
\usepackage[russian]{babel}
\usepackage{graphicx}
\usepackage{array}
\usepackage[a4paper, left=1cm, right=1cm, top=1.27cm, bottom=1.27cm]{geometry}
\usepackage{setspace}
\usepackage{makecell}
\usepackage{caption}

\usepackage{minted}
\usepackage{xcolor}
\usepackage{amsmath}
\usepackage{amssymb}
\geometry{margin=2cm}


\begin{document}

\begin{flushright}
    Пахомов Владимир Юрьевич, 467036
\end{flushright}

\begin{center}
    \textbf{Расчётно-графическая работа №3}
\end{center}

\begin{flushleft}
    \textbf{Условие} \\
    Рассмотрим циклический код над полем $\mathbb{F}_3$ с $n = 13$ и порождающим многочленом $g(X) = (X^3 - X - 1)(X^3 + X^2 -1 )$. Выписать порождающую и проверочную матрицу этого кода и доказать, что он исправляет две ошибки. 
    Декодировать принятое сообщение
    \[ f(X) = 1 - X^3 + X^5 - X^8 - X^9 - X^{12} \]


    \vspace{20pt}
    \textbf{Решение}: \\
    Заменим коэффициенты $-1$ на $2$ (по модулю 3) в многочлене $g(X)$ и раскроем скобки:
    \[ g(X) = (X^3 + 2X + 2)(X^3 + X^2 + 2) \]
    \[ g(X) = X^6 + X^5 + 2X^3 + 2X^4 + 2X^3 + 4X + 2X^3 + 2X^2 + 4 \]
    \[ g(X) = X^6 + X^5 + 2X^4 + 6X^3 + 2X^2 + 4X + 4 \]
    \[ g(X) = X^6 + X^5 + 2X^4 + 2X^2 + X + 1 \]
    $deg(g)=6$. Размерность кода $k = n - deg(g) = 13 - 6 = 7$ \\
    Коэффициенты $g(X): (1, \; 1, \; 2, \; 0, \; 2, \; 1, \; 1)$ \\
    Порождающая матрица $G$: \\
    \begin{table}[h!]
        \centering
        \begin{tabular}{ccccccccccccc}
            1 & 1 & 2 & 0 & 2 & 1 & 1 & 0 & 0 & 0 & 0 & 0 & 0 \\
            0 & 1 & 1 & 2 & 0 & 2 & 1 & 1 & 0 & 0 & 0 & 0 & 0 \\
            0 & 0 & 1 & 1 & 2 & 0 & 2 & 1 & 1 & 0 & 0 & 0 & 0 \\
            0 & 0 & 0 & 1 & 1 & 2 & 0 & 2 & 1 & 1 & 0 & 0 & 0 \\
            0 & 0 & 0 & 0 & 1 & 1 & 2 & 0 & 2 & 1 & 1 & 0 & 0 \\
            0 & 0 & 0 & 0 & 0 & 1 & 1 & 2 & 0 & 2 & 1 & 1 & 0 \\
            0 & 0 & 0 & 0 & 0 & 0 & 1 & 1 & 2 & 0 & 2 & 1 & 1 \\
        \end{tabular}
    \end{table}
    Проверочный многочлен $h(X)$:
    \[ h(X) = \frac{x^{13} - 1}{g(x)} = \frac{x^{13} + 2}{g(x)}\]
    \[ h(X) = X^7 + 2X^6 + 2X^5 + X^2 + X + 2 \]
    Коэффициенты $h(X): (1, \; 2, \; 2, \; 0, \; 0, \; 1, \;1, \; 2$) \\
    Матрица $H$: 
    \begin{table}[h!]
        \centering
        \begin{tabular}{ccccccccccccc}
            1 & 2 & 2 & 0 & 0 & 1 & 1 & 2 & 0 & 0 & 0 & 0 & 0 \\
            0 & 1 & 2 & 2 & 0 & 0 & 1 & 1 & 2 & 0 & 0 & 0 & 0 \\
            0 & 0 & 1 & 2 & 2 & 0 & 0 & 1 & 1 & 2 & 0 & 0 & 0 \\
            0 & 0 & 0 & 1 & 2 & 2 & 0 & 0 & 1 & 1 & 2 & 0 & 0 \\
            0 & 0 & 0 & 0 & 1 & 2 & 2 & 0 & 0 & 1 & 1 & 2 & 0 \\
            0 & 0 & 0 & 0 & 0 & 1 & 2 & 2 & 0 & 0 & 1 & 1 & 2 \\
        \end{tabular}
    \end{table}
    \newpage
    Чтобы исправлять $t$ ошибок, нужно чтобы расстояние $d$ было $d \ge 2t + 1$ \\
    Так как длина кода $n = 13$ и степень порождающего многочлена $6$, то код квадратично-вычетный. Для такого кода выполняется: 
    \[ d^2 \ge n \Rightarrow d^2 \ge 13 \Rightarrow 4 \; (3.6) \]

    Для квадратично-вычетных кодов длины $n \equiv 1 \; \; (mod \; 4)$ минимальное расстояние должно быть нечетным, поэтому берем $d = 5$. Получается так как $d = 5$, то код может исправить $t = \frac{d - 1}{2} = 2$ ошибки. \\
    Теперь декорируем сообщение. Приведем коэффициенты многочлена $f(X)$ в поле $\mathbb{F}_2$:
    \[ f(X) = 1 + 2X^3 + X^5 + 2X^8 + 2X^9 + 2X^{12} \]
    Для нахождения синдрома $S(X)$ найдем остаток от деления $f(X)$ на $g(X)$. Получаем
    \[ S(X) = 2X^5 + X^3 + 2X^2 + 2X + 1 \]
    Теперь найдем вектор ошибок $e(X)$. Остаток от деления $e(X)$ на $g(X)$ равен нашему синдрому $S(X)$:
    \[ e(X) \equiv S(X) \; \; (mod \; g(X)) \]
    В синдроме $S(X)$ слагаемое $2X^5$ имеет степень меньше $g(X)$, то есть если бы это была ошибка, то остаток был равен тому же слагаемому $2X^5$. Вычтем его из синдрома
    \[ S'(X) = X^3 + 2X^2 + 2X + 1 \]
    Этот многочлен является остатком от деления $X^8$ на $g(X)$ (методом перебора). Получается сидром состоит из суммы двух остатков, тогда вектор ошибок $e(X)$:
    \[ e(X) = 2X^5 + X^8 \]

    Для получения исходного слова нужно вычесть ошибку из принятого сообщения:
    \[ c(X) = f(X) - e(X) = f(X) + 2e(X) = \]
    \[ = 1 - X^3 + X^5 - X^8 - X^9 - X^{12} + X^5 + 2X^8 = \]
    \[ = 1 - X^3 - X^5 + X^8 - X^9 - X^{12} \]
    $c(X) = 1 - X^3 - X^5 + X^8 - X^9 - X^{12} \; -$ декодированное сообщение

\end{flushleft}

\end{document}

